\documentclass[11pt]{article}
\usepackage[a4paper,margin=1in]{geometry}
\usepackage{amsmath,amssymb,amsthm,hyperref}
\newtheorem{theorem}{Theorem}
\newcommand{\TTN}{\operatorname{TTN}}
\newcommand{\diag}{\operatorname{diag}}
\title{A counterexample to the tensor-tubal-norm conjecture}
\author{Principia Math In Progress}
\date{22 September 2026}
\begin{document}
\maketitle
\begin{abstract}
The sum of the Euclidean norms of the singular tubes in the standard
t-SVD is not a norm. Two real tensors of size $2\times2\times2$ have
value $\sqrt{5/2}$ each, while their sum has value $2+\sqrt2$.
This also disproves the proposed convex-relaxation assertion.
\end{abstract}

\section*{Task 1: the original conjecture}
Conjecture 2.4.1 in Zhang, Ely, Aeron, Hao and Kilmer
\cite{original}, p.~4, concerns
\[
 F(X)=\TTN(X)=\sum_i\|S(i,i,:)\|_2,
 \qquad X=U*S*V^{\mathsf T}.
\]
Here the t-SVD is calculated by Algorithm 1 of that source: take the
discrete Fourier transform in the third coordinate, compute the usual
matrix SVD of each Fourier slice, and apply the inverse transform.
We use the standard nonincreasing order of the singular values in each
matrix SVD. The conjecture states that $F$ is a norm and is the tightest
convex relaxation of tensor-tubal-rank.

This functional differs from the tensor-nuclear-norm in Theorem 2.4.1
of the same source. The latter sums all singular values of all Fourier
slices. The counterexample below concerns the singular-tube functional
in the conjecture.

\section*{Exact counterexample}
For a tensor with two frontal slices $X_1,X_2$, the Fourier convention
in the source gives
\[
 \widehat X_1=X_1+X_2,\qquad \widehat X_2=X_1-X_2.
\]
Write $s_{ij}=\sigma_i(\widehat X_j)$, where singular values are sorted.
The $i$th diagonal tube of $S$ is
\[
 \left(\frac{s_{i1}+s_{i2}}2,\frac{s_{i1}-s_{i2}}2\right).
 \quad\text{Consequently}\quad
 F(X)=\frac1{\sqrt2}\sum_i\sqrt{s_{i1}^2+s_{i2}^2}.
 \tag{1}
\]
This identity follows directly from the inverse Fourier transform and
the Euclidean norm; it involves no numerical approximation.

\begin{theorem}
The real tensors $A,B\in\mathbb R^{2\times2\times2}$ with frontal slices
\[
 \begin{aligned}
 A_1&=\diag(3/2,0),& A_2&=\diag(1/2,0),\\
 B_1&=\diag(1/2,1),& B_2&=\diag(-1/2,1)
 \end{aligned}
\]
satisfy
\[
 F(A)=F(B)=\sqrt{5/2},\qquad
 F(A+B)=2+\sqrt2>\sqrt{10}=F(A)+F(B).
\]
Thus $F$ is neither a norm nor a convex function.
\end{theorem}

\begin{proof}
Their Fourier slices and those of their sum $C=A+B$ are
\[
\begin{array}{c|cc}
 &\widehat X_1&\widehat X_2\\ \hline
 A&\diag(2,0)&\diag(1,0)\\
 B&\diag(0,2)&\diag(1,0)\\
 C&\diag(2,2)&\diag(2,0).
\end{array}
\]
Both $A$ and $B$ therefore have singular-value lists $(2,0)$ and
$(1,0)$ in the first and second frequency slices. Their diagonal
singular tubes are $(3/2,1/2)$ and $(0,0)$, proving their asserted values.
For $C$ the lists are $(2,2)$ and $(2,0)$, so its diagonal singular
tubes are $(2,0)$ and $(1,1)$. This gives $F(C)=2+\sqrt2$.

All these lists arise from valid real t-SVDs. For $A$ take identity
left and right singular-vector matrices at both frequencies. For $B$
take the coordinate-swap permutation matrix for both left and right
vectors at the first frequency, and identity matrices at the second.
For $C$ take identity matrices at both frequencies. Inverse Fourier
transforms of these real frequency factors produce real tensor
factors; frequencywise orthogonality gives t-orthogonality and the
frequencywise products reconstruct the displayed tensors. The repeated
singular value of $\widehat C_1=2I$ causes no ambiguity in its
singular-value list.

Since both sides are positive, the strict comparison follows by squaring:
$(2+\sqrt2)^2=6+4\sqrt2>10$. Finally, $F$ is positively homogeneous
by (1). Hence
\[
 F\left(\frac{A+B}{2}\right)
 =1+\frac1{\sqrt2}>\sqrt{5/2}
 =\frac{F(A)+F(B)}2,
\]
which violates convexity as well as the triangle inequality.
\end{proof}

A convex envelope must be convex. Thus this functional cannot be the
proposed tightest convex relaxation of tubal rank. Although the source
does not specify a bounding set for that assertion, the obstruction
persists on every norm ball about zero of positive radius: multiply
$A,B$ and their midpoint by the same sufficiently small positive scalar.
This observation addresses the usual bounded-domain interpretation;
the disproof does not depend on the omission of a bounding set.

The source's distinct tensor-nuclear-norm takes values $3,3,6$ on
$A,B,C$, respectively, and is not contradicted. The later conference
paper \cite{conference}, Section 2.4, omits the TTN conjecture and
discusses the different tensor-nuclear-norm. Its omission alone is
not a proof claim resolving the original conjecture.

\section*{Proposed follow-up Tasks 2--5}
These are proposed research tasks, not further solved problems.
\begin{enumerate}
\setcounter{enumi}{1}
\item Determine the convex envelope of $F$ on the Frobenius unit ball
in $\mathbb R^{2\times2\times2}$.
\item Determine the smallest constant $K$ such that
$F(X+Y)\leq K(F(X)+F(Y))$ for all real $2\times2\times2$ tensors.
\item Characterize real linear subspaces of $2\times2\times2$ tensors
on which the restriction of $F$ is a norm.
\item Find the convex envelope of tensor-tubal-rank on the Frobenius
unit ball for real $2\times2\times2$ tensors and compare it with
the normalized tensor-nuclear-norm.
\end{enumerate}

\begin{thebibliography}{9}
\bibitem{original} Z. Zhang, G. Ely, S. Aeron, N. Hao and M. Kilmer,
\emph{Novel Factorization Strategies for Higher Order Tensors:
Implications for Compression and Recovery of Multi-linear Data},
\href{https://arxiv.org/abs/1307.0805v3}{arXiv:1307.0805v3}, 2013.
\bibitem{conference} Z. Zhang, G. Ely, S. Aeron, N. Hao and M. Kilmer,
\emph{Novel methods for multilinear data completion and de-noising
based on tensor-SVD}, CVPR 2014, pp.~3842--3849,
\href{https://doi.org/10.1109/CVPR.2014.485}{doi:10.1109/CVPR.2014.485}.
Official CVF open-access conference copy inspected.
\end{thebibliography}
\end{document}
